\chapter{Nombres premiers, Fermat et chiffrement}
\section*{Exercice 31.1 -- Tester un nombre premier}
[Base] Montrer que 97 est premier.

\section*{Exercice 31.2 -- Factorisation et pgcd}
[Méthode] Décomposer 360 et 252 en facteurs premiers ; en déduire leur pgcd et leur ppcm.

\section*{Exercice 31.3 -- Fermat}
[Méthode] Calculer le reste de $3^{100}$ modulo 7.

\section*{Exercice 31.4 -- Infinité des nombres premiers}
[Preuve] Si $p_1,\ldots,p_r$ est une liste finie de nombres premiers, montrer qu'il existe un nombre premier hors de cette liste.

\section*{Exercice 31.5 -- Chiffrement jouet}
[Application] Pour RSA avec $p=3,q=11,n=33,e=3,d=7$, chiffrer le message $m=2$ et le déchiffrer. Vérifier la relation entre $e,d$ et $(p-1)(q-1)$.

\section*{Exercice 31.6 -- Réciproque de Fermat}
[Critique] Le nombre $341=11\times31$ vérifie $2^{340}\equiv1\pmod{341}$. Le démontrer sans développer la puissance, puis conclure sur un test de primalité fondé sur une seule base.
